\section{Evaluation}

\subsection{Datasets}

To evaluate the proposed method, we conducted experiments across six distinct benchmarks. This selection was designed to perform evaluation across varying domains, structural complexities, reasoning requirements and mitigate potential bias, related to limited domain diversity:
\begin{itemize}
    \item \textbf{Natural Questions}~\cite{kwiatkowski-etal-2019-natural} is a large-scale corpus for open-domain questions developed by Google that consists of over 307K samples where each sample includes a natural language query paired with relevant Wikipedia pages containing the answer spans. The questions originated from real user searches on Google Search Engine. Key distinguishing features compared to other benchmarks include: (1) diversity in question types - NQ encompasses factual, definitional, list-based, comparative, and opinion-oriented queries; (2) complex answer requirements - answers can be short text snippets or long passages requiring deeper reasoning.
    \item \textbf{TriviaQA}~\cite{joshi-etal-2017-triviaqa} is a large-scale benchmark designed for open-domain factoid question answering, featuring over 95K question-answer pairs sourced from Bing search engine. Its distinguishing characteristics include: (1) multi-evidence reasoning - answers often require synthesizing information across multiple documents rather than relying on single-sentence evidence; (2) contextual complexity - diverse types of questions from various domains like history, science, literature. Compared to other datasets like SQuAD or Natural Questions, which focus primarily on extractive question answering within structured contexts, TriviaQA emphasizes multi-hop inference and retrieval-based tasks, making it particularly suitable for evaluating advanced machine reading comprehension systems capable of handling complex queries, requiring broad contextual understanding.
    \item \textbf{HotpotQA}~\cite{yang-etal-2018-hotpotqa} is a crowdsourced question answering dataset built on English Wikipedia, comprising approximately 113K questions. Each question is constructed to require the combination of information from the introductory sections of two Wikipedia articles for answering. The dataset provides two gold paragraphs per question, along with a list of sentences identified as supporting facts necessary to answer the question. HotpotQA includes various reasoning strategies such as bridge questions (involving missing entities), intersection questions (e.g., "what satisfies both property A and property B?") and comparison questions (comparing two entities through a common attribute).
    \item \textbf{2WikiMultihopQA}~\cite{ho-etal-2020-constructing} is a multi-hop question answering dataset that contains complex questions requiring reasoning over multiple Wikipedia paragraphs. Each question is designed to necessitate logical connections across different pieces of information to arrive at the correct answer.
    \item \textbf{MuSiQue}~\cite{10.1162/tacl_a_00475} is a challenging multi-hop QA dataset containing approximately 25K 2–4 hop questions, constructed by composing single-hop questions from five existing single-hop QA datasets. It is designed to feature diverse and complex reasoning paths, requiring models to integrate information from multiple hops to generate correct answers.
    \item \textbf{DiaASQ}~\cite{li-etal-2023-diaasq} consists of user dialogues from a Chinese forum focused on mobile device characteristics. A key feature of this dataset is the inclusion of structured "true statements" that encapsulate the core semantic content of each dialogue. For our evaluation needs we procedurally generate complex multi-hop questions based on these statements.
\end{itemize}

Considering computational and engineering complexity of constructing and traversing large memory graphs, we created manageable yet representative subsets from the original datasets. This step was necessary to enable the iterative experimentation required for tuning multiple retrieval algorithms and LLM configurations within practical resource constraints. The resulting subsets used for knowledge graph construction and evaluation are summarized in Table~\ref{tab:datasets_summ}. Detailed preprocessing steps and datasets statistics are provided in Appendix~\ref{app:datasets}.

\begin{table}[ht!]
\centering
\renewcommand{\arraystretch}{1.3}
\resizebox{.35\textwidth}{!}{%
\begin{tabular}{|c|c|c|}
\hline
\textbf{Dataset} & \textbf{\#qa-pairs} & \textbf{\#documents} \\ \hline \hline
Natural Questions           & 3970                & 2000                \\ \hline
TriviaQA         & 500                 & 4925                \\ \hline
HotpotQA         & 2000                & 3933                \\ \hline
2WikiMultihopQA         & 2000                 & 4596                \\ \hline
MuSiQue         & 1931                & 4185                \\ \hline
DiaASQ           & 4800                & 3483                \\ \hline
\end{tabular}%
}

\caption{Characteristics of prepared datasets for PAI-2 and baselines evaluation}
\label{tab:datasets_summ}
\end{table}


\subsection{Metrics}

Traditional statistical evaluation metrics such as BLEU~\cite{papineni2002bleu}, ROUGE~\cite{lin2004rouge} and Meteor Universal~\cite{denkowski2014meteor} struggle to distinguish syntactically similar, but semantically distinct texts. While semantic methods like BERTScore~\cite{zhang2019bertscore} were introduced to address these limitations, our experiments reveal that BERTScore lacks sufficient differentiability, often failing to capture nuanced distinctions between correct and incorrect answers. Therefore, we adopt \textbf{LLM as a judge}~\cite{zheng2023judging} framework and choose Qwen2.5 7B. The judge evaluates question-answer pairs using a structured prompt containing question, ground truth and generated answer. It labels $1$ for correct answers and $0$ for incorrect ones, and we use accuracy as our main metric. Corresponding LLM-prompts and details are provided in Appendix~\ref{app:judgescore_hyperp}.

To validate reliability of LLM as an evaluative judge, human annotation was conducted for best PAI-2 and HippoRAG 2 experimental setups. The responses generated by the Qwen2.5 7B model were annotated using Overlap-3 metric, with domain experts adhering to the same evaluation criteria as the automated judge. Inter-annotator agreement was quantified using Krippendorff’s $\alpha$~\cite{Krippendorff2011ComputingKA}, yielding a mean $\alpha = 0.935$, which indicates a high degree of assessment reliability. Further evaluation of alignment between the automated judge and human annotators is conducted by computing the Pearson correlation coefficient $r$ between the judge’s scores and the majority vote derived from human annotations. A strong mean correlation of $r = 0.86$ was observed, indicating substantial agreement. Details regarding the annotation procedure are provided in Appendix~\ref{app:human_evaluation}.

Additionally, to measure other characteristics we used several LLM based metrics from RAGAS\footnote{https://docs.ragas.io/en/stable/} library: 
\begin{itemize}
    \item \textbf{Context Relevance} measures whether the retrieved contexts is pertinent to the user query. This is done via two independent LLM-as-a-Judge prompt calls that each rate the relevance on a scale of $0$, $1$ or $2$. The ratings are then converted to a $[0,1]$ scale and averaged to produce the final score. Higher scores indicate that the contexts are more closely aligned with the user's query.
    \item \textbf{Faithfulness} measures how factually consistent an answer is with the retrieved context. It ranges from $0$ to $1$, with higher scores indicating better consistency. An answer is considered faithful if all its claims can be supported by the retrieved context.
    \item \textbf{Groundedness} measures how well an answer is supported or "grounded" by the retrieved contexts. It assesses whether each claim in the answer can be found, either wholly or partially, in the provided contexts: $0$ if answer is not grounded in the context at all; $1$ if answer is partially grounded and $2$ if answer is fully grounded (every statement can be found or inferred from the retrieved context).
\end{itemize}

\subsection{Baselines}

We compare PAI-2 with the following baseline methods:
\begin{itemize}
    \item \textbf{LightRAG}~\cite{guo-etal-2025-lightrag} is a simpler alternative to modern GraphRAG methods that focuses on efficiency. LightRAG is a graph-structured RAG framework that employs a dual-level retrieval system, combining low-level entity retrieval with high-level knowledge discovery. It integrates graph structures with vector representations for efficient retrieval of related entities and their relationships.
    \item \textbf{RAPTOR}~\cite{sarthi2024raptor} is a RAG framework which enhances retrieval via recursive summary and hierarchical clustering into a tree structure. RAPTOR recursively clusters chunks of text based on their vector embeddings and generates text summaries of those clusters, constructing a tree from the bottom up. Nodes clustered together are siblings; a parent node contains the text summary of that cluster.
    \item \textbf{HippoRAG 2}~\cite{gutiérrez2025ragmemorynonparametriccontinual} is a non-parametric continual learning framework that leverages Personalized PageRank algorithm over an open knowledge graph constructed using LLM-extracted triples. It enhances multi-hop reasoning capabilities through sophisticated graph traversal and passage integration mechanisms.
    \item \textbf{PersonalAI 1.0} (PAI-1)~\cite{11479299} is a flexible framework for creating external memory based on a knowledge graph for AI Agents. Building upon AriGraph~\cite{anokhin2024arigraphlearningknowledgegraph} architecture, PAI-1 introduce a novel hybrid graph design that supports both standard edges and two types of hyper-edges, enabling rich and dynamic semantic and temporal representations. Also, it supports diverse retrieval mechanisms, including A*, WaterCircles traversal, BeamSearch and hybrid methods, making it adaptable to different datasets and LLM capacities. 
\end{itemize}

We perform baselines evaluation using Qwen2.5 as LLM backbone for their graph construction and information search algorithms. For specific method we used the following embedder model: LightRAG - \textit{BAAI/bge-m3}\footnote{https://huggingface.co/BAAI/bge-m3}; RAPTOR - \textit{sentence-transformers/multi-qa-mpnet-base-cos-v1}\footnote{https://huggingface.co/sentence-transformers/multi-qa-mpnet-base-cos-v1}; HippoRAG 2 - \textit{facebook/contriever}\footnote{https://huggingface.co/facebook/contriever}; PAI-1 - fusion of \textit{intfloat/multilingual-e5-large}\footnote{https://huggingface.co/intfloat/multilingual-e5-large} and BM25.